\chapter{Language Models in Use}
\label{chap:llm_in_use}

\section{Overview}

The harness supports two providers via an OpenAI-compatible HTTP interface.
The model is selected at runtime with the \texttt{--model} flag; no code
change is needed to switch providers.

\begin{verbatim}
# NVIDIA NIM (Nemotron-3-Ultra-550B)
python run_proofs.py --difficulty Easy --limit 50 \
    --model nvidia/nemotron-3-ultra-550b-a55b

# Groq (gpt-oss-20b)
python run_proofs.py --difficulty Easy --limit 50 \
    --model groq/openai/gpt-oss-20b
\end{verbatim}

A \texttt{groq/} prefix selects Groq and is stripped before the request is
sent, so \texttt{groq/openai/gpt-oss-20b} reaches the API as
\texttt{openai/gpt-oss-20b}. Any other model string goes to whatever
\texttt{QWEN\_BASE\_URL} points at. Provider routing lives entirely in
\texttt{run\_proofs.py}; the harness itself is provider-agnostic.

\section{Provider Details}

\subsection{NVIDIA NIM --- Primary Model}

\begin{itemize}
  \item \textbf{Model:} \texttt{nvidia/nemotron-3-ultra-550b-a55b} (550 billion parameters)
  \item \textbf{API base:} \texttt{https://integrate.api.nvidia.com/v1}
  \item \textbf{Credential:} \texttt{QWEN\_API\_KEY} in \texttt{C:\textbackslash lean\textbackslash .env}
  \item \textbf{Completion budget:} \texttt{max\_tokens=8192} per call
\end{itemize}

The credential variable is named \texttt{QWEN\_API\_KEY} for historical
reasons --- the harness originally targeted Qwen 2.5 on DashScope, and the
variable names (\texttt{QWEN\_API\_KEY}, \texttt{QWEN\_BASE\_URL},
\texttt{QWEN\_MODEL}) were never renamed when the provider changed. They are
provider-neutral in effect: whatever base URL and key they carry is what the
harness talks to.

The budget is generous because Nemotron is a reasoning model. It emits on the
order of 1700 reasoning tokens on set-theory goals before committing to a
tactic; earlier settings of 150 and then 2048 both truncated it mid-sentence,
and the truncated fragment was sent to Lean as though it were a tactic.

\subsection{Groq --- Second Model}

\begin{itemize}
  \item \textbf{Model:} \texttt{openai/gpt-oss-20b} (accessed via the \texttt{groq/} prefix)
  \item \textbf{API base:} \texttt{https://api.groq.com/openai/v1}
  \item \textbf{Credential:} \texttt{GROQ\_API\_KEY} in \texttt{C:\textbackslash lean\textbackslash .env}
  \item \textbf{Limit:} 8000 tokens per minute on the free tier; the harness
        caps \texttt{max\_tokens=1024} to stay under it.
\end{itemize}

The cap matters more than it looks. Groq counts \texttt{max\_tokens} against
the per-minute ceiling \emph{before} the request runs, so a budget that
overshoots is rejected outright rather than truncated:

\begin{verbatim}
Error code: 413 - Request too large for model `openai/gpt-oss-20b` ...
tokens per minute (TPM): Limit 8000, Requested 8381
\end{verbatim}

Once history accumulates the prompt alone runs to roughly 6--7k tokens, which
leaves about 1k for the completion. Carrying the NVIDIA budget of 8192 across
to Groq produced this 413 on 41 of 42 LLM failures in one run, and no amount
of retrying would have helped --- the request was structurally too large. The
budget is therefore resolved per provider at harness construction, and can be
overridden for a single run with \texttt{--max-tokens}.

\section{Shared Configuration}

Both providers receive the same system prompt every step:

\begin{verbatim}
You are proving a Lean 4 theorem one tactic at a time.
You will be shown the current goal state and a history of what you tried.
Respond with EXACTLY ONE Lean 4 tactic and nothing else -- no explanation,
no code fences, no markdown. Just the tactic text, e.g.:
    rw [Nat.add_comm]

If the previous step returned an error, read it carefully and adjust.
If the goal did not change, your tactic had no effect -- try something different.
\end{verbatim}

The full tactic history and current proof state are included in the user turn.
Up to 3 retries are issued if the extracted tactic is empty, each appending a
short nudge to the user message.

\subsection{Tactic extraction}

Both models ignore the ``exactly one tactic'' instruction often enough that the
completion cannot be sent to Lean as-is. \texttt{\_extract\_tactic} recovers the
tactic by scanning the non-empty lines \emph{bottom-up} --- a reasoning model
states its answer last --- after discarding code-fence markers and inline
backticks. A line is accepted when it either opens with a known tactic keyword
followed by call syntax, or is short and free of sentence punctuation. If no
line qualifies, the last non-empty line is returned, so a single-line response
always survives.

The keyword test alone is not sufficient, because the models narrate
\emph{about} tactics constantly. \texttt{simp will rewrite the goal} and
\texttt{exact because the term has type} both open with a tactic keyword and
are both prose; what follows the keyword decides it. The raw completion is
logged verbatim alongside the extracted string, so a bad tactic can be
attributed to the model or to the extractor rather than guessed at.

\section{Reliability and Circuit Breaker}

Neither provider guarantees uptime, and the two fail differently. The harness
therefore counts two kinds of failure separately, because conflating them is
expensive:

\begin{center}
\begin{tabular}{@{}p{0.20\linewidth}p{0.40\linewidth}p{0.34\linewidth}@{}}
\hline
\textbf{Kind} & \textbf{Trigger} & \textbf{Response} \\
\hline
\textbf{Model failure} & The endpoint answered with nothing usable: empty
completion after retries, 404, 401, bad model name & Fail fast.
\texttt{MAX\_CONSECUTIVE\_LLM\_FAILURES = 3} consecutive, then
\texttt{FAILED\_LLM\_ERROR} \\
\textbf{Connectivity failure} & The endpoint was unreachable: connection
error, read timeout, 5xx, 429 & Back off 30/60/120/240/300\,s and retry.
\texttt{MAX\_CONSECUTIVE\_CONNECTIVITY\_FAILURES = 10} (roughly 32 minutes),
then \texttt{FAILED\_CONNECTIVITY\_OUTAGE} \\
\hline
\end{tabular}
\end{center}

Only a successful call resets a counter; a model failure does not hand the
connectivity budget a free reset. A second, coarser breaker sits above these
in \texttt{BatchRunner}: five consecutive proofs ending
\texttt{FAILED\_LLM\_ERROR} abort the whole batch, leaving the remaining
theorems untouched and re-runnable rather than burned against a dead endpoint.

Timeouts are enforced on both sides of the loop:
\texttt{LLM\_CALL\_TIMEOUT\_SECONDS = 600} bounds a single completion request
(one call once hung for 22.3 hours without it), while the Lean REPL gets
\texttt{LEAN\_ELABORATION\_TIMEOUT\_SECONDS = 900} for the initial
elaboration --- which re-imports Mathlib inside the command, measured at
roughly 230\,s --- and \texttt{LEAN\_TACTIC\_TIMEOUT\_SECONDS = 120} for each
tactic step.

\section{What the two Easy-50 runs cost}

\begin{center}
\begin{tabular}{@{}p{0.26\linewidth}p{0.16\linewidth}p{0.16\linewidth}p{0.36\linewidth}@{}}
\hline
\textbf{Run} & \textbf{Headline} & \textbf{Reachable} & \textbf{Lost, and to what} \\
\hline
NIM, Nemotron-550B & 7/50 = 14.0\% & 7/22 = 31.8\% & 28: twenty-four to a
multi-day 404/429/503 outage, four to Lean elaboration \\
Groq, gpt-oss-20b & 8/50 = 16.0\% & 8/28 = 28.6\% & 22: fifteen connectivity
outages, three model failures, four to Lean elaboration \\
\hline
\end{tabular}
\end{center}

The headline figure divides successes by theorems the model never saw, so it
measures provider uptime as much as capability. The reachable-only figure
counts only theorems that reached a working model --- those proved plus those
that exhausted the step budget --- and is the number this study reports. Both
land near the SorryDB agentic baseline of roughly 20--30\% on Easy, on samples
of 22 and 28 respectively, so the confidence intervals are wide and the two
models are not yet distinguishable from one another.

The Groq run took roughly 42 hours, most of it spent in connectivity backoff
rather than computation. That is the retry layer working as designed: it keeps
theorems alive at the cost of wall-clock, where the alternative is discarding
them.
